\chapter{出版说明}

当前，中学教学改革已经深入到课程设置和教材改革领域。
我校数学教材的改革，以发展学生的数学思维为目标，以不改变现行教学大纲规定的教学闪容为前提，试图通过对知识结构及其展开方式的统盘考虑，实现整体优化。经多年反复探索、实验，编成了这套尝试融教材与教法、学法于一体的《北京四中高中教学讲义》。

这套讲义的产生可以上溯到 1982 年。从那时起，为了发展学生智能，提高数学素养，我校部分同志就开始对高中数学教学进行以教材改革为龙头，以学法教育为重点的“整体优化实验研究”。正是在这项研究的基础上，逐步形成了这套讲义编写的特色和风格。这就是：
\begin{enumerate}
    \item 为形成学生良好的认知结构，讲义的知识结构力求脉络分明，使学生能从整体上理解教材。
    \item 为了提高学生的数学素养，本讲义把数学思想的阐述放到了重要位置。数学思想既包含对数学知识点（概念、定理、公式、法则和方法）的本质认识，也包含对问题解决的数学基本观点。它是数学中的精华，对形成和发展学生的数学能力具有特别重要的意义。为此，讲义注重展现思维过程（概念、法则被概括的过程，教学关系被抽象的过程，解题思路探索形成的过程）。在过程中认识知识点的本质，在过程中总结思维规律，在过程中揭示数学思想的指导作用。力图使学生能深刻领悟教材。
    \item “再创造，再发现”在数学学习中对培养创造维能力至关重要，为引导学生积极参与“发现”，讲义在设计上做了某些尝试。
    \item 例题和习题的选配，力求典型、适量、成龙配套。习题分为 A 组（基本题）、 B 组（提高题）和 C 组（研究题）。教师可根据学生不同的学习水平适当选用。
    \item 教材是学生学习的依据。应有利于培养自学能力，本书注重启迪学法，并在书末附有全部习题的答案或提示，以供学习时参考。
 \end{enumerate}   

这套讲义在研究、试教和成书的过程中，始终得到了北京市和西城区教育部门有关领导的关怀和帮助，得到了北京师范大学数学系钟善基教授、曹才翰教授的热情指导，清华附中的瞿宁远老师也积极参与了我们的实验研究，并对这套教材做出了贡献，在此一并致以诚挚的谢意。

在编写过程中，北京四中数学组的教师们积极参加研讨，对他们的热情支持表示感谢。

这套讲义包括六册：高中代数第一、二、三册，三角、立体几何、解析几何各一册。

编写适应素质教育的教材，对我们来说是个尝试。由于水平所限，书中不当之处在所难免，诚恳希望专家、同行和同学们提出宝贵意见。

\begin{flushright}
    北京四中教学处\\
    1994年11月
\end{flushright}

\chapter{绪论}

\subsection*{立体几何研究的对象}
几何学是研究几何图形的性质及其应用的一门学科．

在初中我们学过平面几何． 在那里研究的是平面图形（即所有的点都在同一平面内的图形）的性质（即形状、大小和位置关系）和应用．可是在实际问题中，大量的几何图形并不是平面图形，如房屋、机器、堤坝、桌椅、橡皮等等，它们都可看成几何体．它们是空间图形．

空间图形由空间的点，线，面构成． 也可以看做是空间点的集合． 平面图形是空间图形的一部分． 有些空间图形就是由平面图形围成的． 例如六个正方形围成了一个正方体．有些空间图形不是由平面图形围成，例如球．

立体几何研究的对象就是空间图形．

\subsection*{立体几何课的内容}
空间图形是千姿百态的，从哪儿入手研究它们呢？我们不能今天研究桌子，明天研究凳子，……那样不是科学的研究方法．

我们要首先研究构成空间图形的基本元素：点，线，面．主要是点，直线和平面，弄清了它们的性质及其相互关
系，我们就可驾驭简单的图形，进而可以研究较为复杂的图形．

点，直线和平面之间的位置关系可分为六种：（1）点与点；（2）点与直线；（3）点与平面；（4）直线与直线；（5）直线与平面；（6）平面与平面．

研究上述这六种关系，就构成了立体儿何最基本的内容，即本书第一章的内容。在这个基础上再研究一些简单体的几何性质．就是本书第二章的内容。在本书第三章，研究一类特殊的几何体——正多面体。

\subsection*{立体几何的研究方法}
研究立体几何的方法与研究平面几何的方法类似，不是依靠观察和试验，而是依据公理，运用逻辑推理方法，即：
\begin{enumerate}
    \item 先给出一些最初的、不予定义的基本概念；
    \item 再逐步建立一些能较直观地反映这些基本概念性质的公理；
    \item 然后用这些基本概念去定义新的概念；
    \item 再根据基本概念、公理以及定义的新概念去证明定理，并在此基础上再证明新的定理。
\end{enumerate}

需要指出的是，作为中学数学课程的平面几何与立体几何，限于读者的认识水平和接受能力，不可能也没有必要完全依照严格的公理体系处理，为了便于接受，保证读者能集中精力学好主要内容，我们省略了某些定理的繁复的证明，有些定理直接以公理的形式出现，中学的几何体系只是粗略的公理体系．因为平面图形是空间图形的一部分，是在同一平面内的特殊的空间图形．所以平面几何的基本概念和公理体系，也就自然地包含在更一般的立体几何的基本概念和公理
体系之中． 平面几何的所有定理在立体几何中解决平面问题时都可以直接应用．

\subsection*{学习立体几何需要注意的几个问题}
由研究平面图形到研究空间图形，开始阶段有些同学会感到困难．主要是画图、识图这两个环节，画图是要把空间的立体图形画在一个平面（如纸面或黑板）上；识图就是要根据画在平面上的“立体”图形想象出原来空间图形的真实形状．再就是研究空间图形，由于对象更加复杂，对抽象思维能力和逻辑思维能力的要求更高了，知道了这些特点，我们
就可以采取相应的方法，顺利地进入立体几何这个新的领域，“入了门”就会感到越学越有兴趣了．

\begin{enumerate}
    \item 为了培养自己的空间想象力，为了能“画图”就要结合对空间图形的观察，想象这些空间图形画在纸上应是什么模样；同时要掌握画直观图的规则，掌握实线、虚线的使用规则．从画简单的图形（如直线和平面的各种位置关系）、简单的几何体（如正方体）画起．对照实物或模型观察各种线线、线面、面面关系，然后过渡到根据直观图来想象这些线线、线面、面面关系（这就是识图的过程），最后在没有模型摆在面前，甚至不看直观图，头脑中都能清晰地想象出用语言文字描述的图形．这就是提高空间想象力的过程，也是提高抽象思维能力的过程．
    \item 要重视看起来简单的那些基本概念、公理和定理，不仅要理解它们，还要熟练地记忆它们．要能用图形、符号、语言三种形式来准确地表达它们．对条件和结论要一清二楚，对定理要会证明，并在应用中进一步领会和加深理解．
    \item 要注意与平面几何的联系和区别．立体几何与平面儿何是同一个公理体系，所用的方法都是逻辑推理的方法，所不同的从表面上看可以说是“部分”和“整体”的问题．在平面几何中的正确命题在立体几何中有的就不正确了，有的虽然仍然正确但意义更丰富了，有的只能类比地推广到立体几何中去……这些我们还将在第一章的小结中论及，现在提出这些只是提醒同学们在学习过程中，若能注意它们的联系和区别，及时进行比较和总结，对学好立体几何是非常有益的．
\end{enumerate}




